% ==================================================================
% id:       cond-010
% topic:    הסתברות מותנית ואי-תלות
% subtopic: נוסחת ההסתברות השלמה ונוסחת בייס
% solution: full
% ==================================================================
חשוד משקר בהסתברות 0.8. מעבירים אותו בדיקת פוליגרף. ידוע ש-90\% מהמשקרים נכשלים בבדיקה, ורק 5\% מדוברי האמת נכשלים בה. החשוד נכשל. מה ההסתברות שהוא משקר?
\begin{solbox}
נסמן $L$ = ״משקר״, $F$ = ״נכשל״. נתון $\PP(L)=0.8$, $\PP(F\mid L)=0.9$, $\PP(F\mid L^C)=0.05$. לפי נוסחת ההסתברות השלמה
\[
\PP(F)=0.9\cdot0.8+0.05\cdot0.2=0.72+0.01=0.73,
\]
ולפי נוסחת בייס
\[
\PP(L\mid F)=\frac{\PP(F\mid L)\PP(L)}{\PP(F)}=\frac{0.72}{0.73}\approx\ans{0.986}.
\]
\end{solbox}
